\documentclass[french,a4paper]{article}

\usepackage[utf8]{inputenc}
\usepackage[T1]{fontenc}
\usepackage{lmodern}
\usepackage{geometry}
\usepackage{graphicx}
\usepackage{babel}
\usepackage{amsmath}
\usepackage{amsthm}
\usepackage{amssymb}
\usepackage{hyperref}
\usepackage{stmaryrd}
\usepackage{enumitem}
\usepackage{tcolorbox}
\usepackage{dsfont}
\usepackage{multirow}
\usepackage{etoc}
\usepackage{titlesec}
\usepackage{esint}
\usepackage{cancel}
\usepackage{mathdots}

\setlist[itemize]{label=$\bullet$}


\renewcommand{\P}{\mathbb{P}}
\newcommand{\N}{\mathbb{N}}
\newcommand{\Z}{\mathbb{Z}}
\newcommand{\Q}{\mathbb{Q}}
\newcommand{\R}{\mathbb{R}}
\newcommand{\C}{\mathbb{C}}
\newcommand{\K}{\mathbb{K}}
\renewcommand{\L}{\mathbb{L}}
\newcommand{\U}{\mathbb{U}}
\newcommand{\st}{^*}
\newcommand{\tq}{\middle|}
\newcommand{\inv}{^{-1}}
\newcommand{\E}{\mathbb{E}}
\newcommand{\F}{\mathbb{F}}
\newcommand{\V}{\mathbb{V}}
\renewcommand{\ker}{\text{Ker}}
\newcommand{\im}{\text{Im}}
\newcommand{\card}{\text{Card}}
\newcommand{\Tr}{\text{Tr}}

\theoremstyle{plain}
\newtheorem{thm}{Théorème}
\newtheorem*{ind}{Indication}

\theoremstyle{definition}
\newtheorem{dfn}{Definition}

\theoremstyle{remark}
\newtheorem*{rmq}{Remarque}
\newtheorem*{app}{Application}

\newtcolorbox[auto counter]{ex}[2][]{colback=white, colframe=green!50!white, coltitle=black, fonttitle=\bfseries, title={Exercice~\thetcbcounter~: #2}, after title={\hfill #1}}

\title{TD Polynômes et fractions rationnelles}

\begin{document}

\maketitle

Dans toute cette feuille, sauf mention explicite du contraire, $\K$ désigne un surcorps de $\Q$.

\section{Algèbre des polynômes}

\begin{ex}[Moulin.Pol.01]{$\star$ Une équation polynomiale}
	Déterminer tous les polynômes $P\in\C[X]$ tels que $P(X^2)=P(X)^2$.
\end{ex}
\begin{proof}
	Les polynômes constants égaux à $0$ et $1$ conviennent évidemment, et aucun autre polynôme constant ne convient.\\ On se donne alors un tel polynôme $P$ non constant. On remarque que si $P$ possède une racine $z$, alors pour tout $n\in\N$, $z^{2^n}$ est racine de $P$. Donc $P$ ne peut pas posséder de racine non nulle de module différent de $1$, sinon il en possèderait une infinité et serait donc nul. Les racines de $P$ sont donc nulles ou de module $1$. Mais comme $P$ possède un nombre fini de racines, si $z$ est une racine de $P$ de module $1$, alors il existe $p<n$ tel que $z^{2^p}=z^{2^n}$. Donc $z$ est une racine de l'unité. Mais de la même manière, si $e^{i\theta}$ est racine de $P$, alors $e^{i\theta/2}$ est racine de $P$. Pour ne pas créer une infinité de racines en divisant les arguments par $2$ successivement, $1$ est la seule racine possible de module $1$. Mais alors, $-1$ serait racine si $1$ l'était, ce qui serait alors absurde car la seule racine de module $1$ potentielle était $1$. Donc la seule racine possible pour $P$ est $0$. Et comme $P$ est non constant, on déduit de $P(X^2)=P(X)^2$ le fait que son coefficient dominant vaut $1$. \\
	Réciproquement, les polynômes $0$, $1$ et $X^n$ pour $n\in\N\st$ conviennent bien.
\end{proof}

\begin{ex}[Moulin.Pol.02]{Une équation polynomiale}
	Déterminer les $P\in\R[X]$ tels que $P(0)=0$ et $P(X^2+1)=P(X)^2+1$.
\end{ex}
\begin{proof}
	Montrer que le polynôme $X$ est la seul solution (c'est évidemment une solution). \\
	Pour cela, remarquer que $P(u_n)=u_n$ pour tout $n\in\N$ où $(u_n)_{n\in\N}$ est la suite définie par : $$\begin{cases}
		u_0=0 \\
		\forall n\in\N,\ u_{n+1}=u_n^2+1
	\end{cases}$$ Cette suite est strictement croissante (utiliser un trinôme du second degré) donc injective, si bien que $P$ et $X$ sont égaux puisqu'ils coexistent en une infinité de points.
\end{proof}

\begin{ex}[Moulin.Pol.03]{Polynômes positifs}
	Soit $P\in\R[X]$ tel que $\forall x\in\R,\ P(x)\geqslant 0$. Montrer qu'il existe $(A,B)\in\R[X]^2$ tel que $P=A^2+B^2$. Indication : on pourra commencer par remarquer que l'ensemble des polynômes qui sont somme de deux carrés est stable par multiplication.
\end{ex}
\begin{proof}
	Pour commencer, on peut remarquer comme le suggère l'indication qu'un produit de deux sommes de deux carrés est une somme de deux carrés. En effet, cela est valable dès que l'anneau dans lequel on travaille est commutatif : $$(a^2+b^2)(c^2+d^2)=(ac-bd)^2+(ad+bc)^2$$ (on peut interpréter ceci en termes de modules de nombres complexes si on le souhaite, mais cette factorisation peut servir). Par une récurrence immédiate, cela est valable pour un nombre fini de facteurs. \\
	Ensuite, on décompose $P$ en facteurs irréductibles sur $\R[X]$ (on suppose $P$ non constant sinon le résultat est trivial). Par un argument de limite en $+\infty$, le coefficient dominant de $P$ est (strictement) positif. Les multiplicités dans les facteurs de degré $1$ sont nécessairement paires car sinon le polynôme changerait de signe en ces racines. Enfin, pour les trinômes du second degré irréductibles, on les met sous forme canonique. On applique ensuite le résultat préliminaire et le tour est joué.
\end{proof}

\begin{ex}[Moulin.Pol.04]{$\star\star$}
	Soit $P\in\R[X]$ tel que $\forall x\in\R_+,\ P(x)>0$. Montrer qu'il existe $N\in\N$ tel que $(1+X)^NP$ soit à coefficients positifs.
\end{ex}
\begin{proof}
	Toutes les racines réelles de $P$ sont donc strictement positives et le coefficient dominant de $P$ est strictement positif. En décomposant $P$ en facteurs irréductibles sur $\R$, on peut donc écrire : $$P=\lambda\prod_{k=1}^{r}(X+\alpha_k)\prod_{l=1}^{s}(X^2+\beta_lX+\gamma_l)$$ avec $\lambda>0$, les $\alpha_k$ strictement positifs et $\beta_l^2-4\gamma_l<0$. \\ Pour qu'un polynôme réel soit à coefficients positifs, il suffit qu'il soit le produit de polynômes à coefficients positifs. Les facteurs irréductibles de degré $1$ le sont déjà, donc il suffit de voir l'effet de $(1+X)^N$ sur un trinôme du second degré irréductible. Après utilisation d'un binôme de Newton, on trouve que le coefficient de $X^k$ dans $(1+X)^n(X^2+\beta X+\gamma)$ est : $$x_{k,N}=\gamma\binom{N}{k}+\beta\binom{N}{k-1}+\binom{N}{k-2}$$ On traite d'abord le cas où $k\in\llbracket 2,N\rrbracket$. On a alors : $$x_{k,N}=\underbrace{\frac{N!}{(N-k+2)!k!}}_{\geqslant 0}\times\underbrace{(\gamma(N-k+2)(N-k+1)+\beta k(N-k+2)+k(k-1))}_{=:P_N(k)}$$ $P_N(k)$ est un polynôme du second degré, on va essayer de faire en sorte qu'il soit à valeurs positives. Sous forme développée, $P_n$ s'écrit : $$P_N=(\gamma-\beta+1)X^2+(\beta(N+2)-1-\gamma(2N+3))X+\gamma(N+2)(N+1)$$ Le discriminant de de $P_n$ vaut alors après simplification : $$\Delta_N=(\beta^2-4\gamma)N^2+O(N)$$ Comme $\beta^2-4\gamma<0$ par hypothèse, $\Delta_N$ tend vers $-\infty$ quand $N$ tend vers $+\infty$, donc est strictement négatif à partir d'un certain rang $N_0$. Or, le coefficient dominant de $P_N$ est exactement le trinôme $X^2+\beta X+\gamma$ évalué en $-1$, et sa valeur est strictement positive par hypothèse. Donc pour tous les $N\geqslant N_0$, $P_N$ est à valeurs strictement positives et les $P_N(k)$ sont donc positifs. \\ Enfin, traitons les quelques cas particuliers restants : \begin{itemize}
		\item $k=0$ : le terme constant vaut $\gamma>0$ ;
		\item $k=1$ : le terme de degré $1$ vaut $N\gamma+\beta>0$ pour $N$ supérieur à un certain $N_1$ ;
		\item $k=N+1$ : le terme de degré $N+1$ vaut $N+\beta>0$ pour un $N$ supérieur à un certain $N_2$ ;
		\item $k=N+2$ : le terme de degré $N+2$ vaut $1>0$.
	\end{itemize}
	Ainsi, pour $N= \max(N_0,N_1,N_2)$, $(1+X)^N(X^2+\beta X+\gamma)$ est à coefficients positifs. \\ On conclut en prenant $N=N_1+\cdots+N_s$ où $N_l$ est un entier qui convient pour le trinôme irréductible $X^2+\beta_l X+\gamma_l$.
\end{proof}

\begin{ex}[Moulin.Pol.05]{$\Delta$}
	\begin{enumerate}
		\item $\circ$ Déterminer les polynômes $P\in\C[X]$ tels que $P(\C)=\C$.
		\item Déterminer les polynômes $P\in\C[X]$ tels que $P(\R)=\R$.
		\item $\star\star$ Déterminer les polynômes $P\in\C[X]$ tels que $P(\Q)=\Q$.
	\end{enumerate}
\end{ex}
\begin{proof}
	Exercice très classique !
	\begin{enumerate}
		\item Ce sont les polynômes non constants (en utilisant le théorème de d'Alembert-Gauss).
		\item $P$ et $\overline{P}$ coïncident sur $\R$ qui est une partie infinie de $\C$ donc $P$ est à coefficients réels (on pourrait le faire avec une interpolation de Lagrange mais c'est plus moche). Ensuite, ce sont les polynômes de degré impair (utiliser le TVI et les limites en $\pm\infty$). Les polynômes de degré pair sont en effet toujours majorés ou minorés.
		\item Ce sont exactement les polynômes de degré $1$ à coefficients rationnels. Déjà, un tel polynôme est à coefficients rationnels par interpolation de Lagrange. Sans perte de généralité, on suppose $P$ à coefficients entiers, quitte à la multiplier par un nombre entier non nul, et on le prend non constant, de degré $n$ (sinon le résultat est trivialement faux). On regarde ensuite les antécédents de $1/p$ pour $p$ premier : $$\frac{1}{p}=P\left(\frac{r}{s}\right)=\sum_{k=0}^{n}a_k\frac{r^k}{s^k}$$ avec $r\wedge s=1$. On peut alors réécrire ceci en : $$s^n=p\sum_{k=0}^{n}a_kr^ks^{n-k}$$ Donc $p\mid s^n$ et donc $p\mid s$ car $p$ est premier. On écrit alors $s=pd$ puis on simplifie par $p$ : $$p^{n-1}d^n=\sum_{k=0}^{n}a_kr^kp^{n-k}d^{n-k}$$ Par l'absurde, on suppose $n\geqslant 2$. En réduisant modulo $p$, on trouve que $p$ divise $r^na_n$, mais comme $p$ divise $s$ et $r\wedge s=1$, $p\wedge r^n=1$ donc $p$ divise $a_n$. Mais alors, on vient de prouver que $a_n$ est divisible par tous les nombres premiers, donc que $a_n=0$ ce qui est absurde. Donc $n=1$. Et réciproquement, les polynômes de degré $1$ à coefficients rationnels conviennent bien.
	\end{enumerate}
\end{proof}

\begin{ex}[Moulin.Pol.06]{$\star$ Que des $1$ en base $10$}
	On note $A$ l'ensemble des entiers naturels ne s'écrivant en base $10$ qu'avec des $1$. Déterminer les polynômes $P\in\R[X]$ stabilisant $A$.
\end{ex}
\begin{proof}
	Pour $n\in\N\st$, on pose $a_n=(10^n-1)/9$. Alors, $A=\left\{a_n\mid n\in\N\st\right\}$. Soit $P$ de degré $d$ ($d\geqslant 0$ car le polynôme nul ne convient pas) tel que $P(A)\subset A$. Alors, pour tout $n\in\N$, il existe $m_{n}\in\N\st$ tel que $P(u_n)=u_{m_n}$. Comme la suite $(u_n)$ tend vers $+\infty$, en notant $\lambda$ le coefficient dominant de $P$, on trouve en $+\infty$ : $$\frac{10^{m_n}}{9}\sim u_{m_n}\sim \lambda u_n^d\sim\lambda\frac{10^{nd}}{9^d}$$ On en déduit : $$10^{m_n-nd}\sim\frac{\lambda}{9^{d-1}}$$ On suppose dans un premier temps $d\geqslant 1$. En passant au logarithme dans les limites :  $$m_n-nd\sim\log_{10}\left(\frac{\lambda}{9^{d-1}}\right)$$ Comme $m_n-nd$ est un entier et qu'une suite d'entiers qui converge stationne vers sa limite, on en déduit qu'à partir d'un certain rang : $$m_n=nd+\log_{10}\left(\frac{\lambda}{9^{d-1}}\right)$$ et que $k=\displaystyle\log_{10}\left(\frac{\lambda}{9^{d-1}}\right)\in\Z$.
	Mais alors, comme les $x_n=10^n$ forment une partie infinie de $\R$, on en déduit : $$P\left(\frac{X-1}{9}\right)=\frac{10^kX^d-1}{9}$$ En effectuant le changement de variable $T=\displaystyle\frac{X-1}{9}$, on trouve : $$P(T)=\frac{10^k(9T+1)^d-1}{9}$$ Et en regardant pour $T=1$, on trouve que l'on doit avoir $k\geqslant 1-d$. \\ On vérifie aisément que les $$\frac{10^k(9X+1)^d-1}{9}$$ conviennent pour $d\in\N$ et $k\geqslant 1-d$.
\end{proof}

\begin{ex}[Morlot]{Un polynôme toujours positif}
	Soit $P$ un polynôme de degré $n$ positif sur tout $\R$. Montrer que $Q=P+P'+\ldots+P^{(n)}$ est positif sur tout $\R$. 
\end{ex}
\begin{proof}
	Dériver $e^{-x}Q(x)$ en remarquant que $Q'-Q=P$. Utiliser un argument de limite en $+\infty$ pour conclure.
\end{proof}

\section{Relations coefficients-racines}

\begin{ex}[Moulin.Pol.07]{Système d'équations \textit{via} les polynômes}
	Résoudre le système d'équations : $$\begin{cases}
		a+b+c=2 \\
		a^2+b^2+c^2=6 \\
		a^3+b^3+c^3=8
	\end{cases}$$
\end{ex}

\begin{ex}[Moulin.Pol.08]{CNS pour qu'une racine soit somme des autres}
	Donner une condition nécessaire et suffisante sur $(p,q,r)\in\C^3$ pour que l'une des racines du polynôme $X^3+pX^2+qX+r$ soit la somme des deux autres.
\end{ex}

\begin{ex}[Moulin.Pol.09]{$\star$ Triangle équilatéral}
	Donner une condition nécessaire et suffisante sur $(p,q,r)\in\C^3$ pour que les racines du polynôme $X^3+pX^2+qX+r$ forment un triangle équilatéral (éventuellement réduit à un point).
\end{ex}

\begin{ex}[Moulin.Pol.10]{$\star$ Parallélogramme}
	Donner une condition nécessaire et suffisante sur $(p,q,r,s)\in\C^4$ pour que les racines du polynôme $X^4+pX^3+qX^2+rX+s$ forment un parallélogramme (éventuellement plat).
\end{ex}

\begin{ex}[Moulin.Pol.11]{$\star\star$ Inégalité de Samuelson}
	Soit $n\in\N$ avec $n\geqslant 2$ et $P=\displaystyle X^n+\sum_{k=1}^{n}a_kX^{n-k}\in\R[X]$. On suppose que toutes les racines de $P$ sont réelles. Montrer qu'elles appartiennent à l'intervalle $$\left[-\frac{a_1}{n}-\frac{n-1}{n}\sqrt{a_1^2-\frac{2n}{n-1}a_2},-\frac{a_1}{n}+\frac{n-1}{n}\sqrt{a_1^2-\frac{2n}{n-1}a_2}\right]$$
\end{ex}
\begin{proof}
	Notons $x_1,\ldots,x_n$ ses racines comptées avec multiplicité. On pose : $$s_1=\sum_{i=1}^{n}x_i\ \ \text{et}\ \ s_2=\sum_{i\ne j}x_ix_j$$ D'après les formules de Viète : $$s_1=-a_1\ \ \text{et}\ \ s_2=a_2$$ On en déduit : $$\sum_{i=1}^{n}x_i^2=\left(\sum_{i=1}^{n}x_i\right)^2-2\sum_{i\ne j}x_ix_j=a_1^2-2a_2$$ On fixe une racine $x_1$ (sans perte de généralité quitte à permuter les $x_i$) que l'on isole puis on applique l'inégalité de Cauchy-Schwarz (car toutes les racines sont réelles !) : $$(-a_1-x_1)^2=\left(\sum_{i=2}^{n}x_i\right)\leqslant(n-1)\sum_{i=2}^{n}x_i^2=(n-1)(a_1^2-2a_2-x_1^2)$$
	Donc : $$a_1^2+2a_1x_1+x_1^2-(n-1)(a_1^2-2a_2-x_1^2)=nx_1^2+2a_1x_1+((2-n)a_1^2+2(n-1)a_2)\leqslant 0$$ On s'intéresse aux racines du polynôme : $$Q=nX^2+2a_1X+((2-n)a_1^2+2(n-1)a_2)$$ Le discriminant vaut :
	\begin{align*}
		\Delta&=(2a_1)^2-4n((2-n)a_1^2+2(n-1)a_2)\\
		&=4a_1^2-4n((2-n)a_1^2+2(n-1)a_2)\\
		&=(4n^2-8n+4)a_1^2-8n(n-1)a_2\\
		&=[2(n-1)]^2\left(a_1^2-\frac{2n}{n-1}a_2\right)
	\end{align*}
	Les racines de $Q$ sont donc : $$-\frac{a_1}{n}\mp\frac{n-1}{n}\sqrt{a_1^2-\frac{2n}{n-1}a_2}$$ Or, $Q(x_1)\leqslant 0$ et $Q$ est un polynôme du second degré de coefficient dominant positif, donc $x_1$ est situé entre les racines de $Q$. Cela étant valable pour toutes les racines de $P$, on a bien le résultat souhaité.
\end{proof}

\section{Racines complexes}

\begin{ex}[Moulin.Pol.12]{$\Delta$ CNS pour être scindé sur $\R$}
	Soit $P\in\R[X]$ unitaire de degré $n\geqslant1$. Alors, $P$ est scindé sur $\R$ si, et seulement si, $$\forall z\in\C,\ \lvert P(z)\rvert\geqslant\lvert\im(z)\rvert^n$$
\end{ex}
\begin{proof}
	On procède par double implication :
	\begin{itemize}
		\item Sens direct : si $P$ est scindé sur $\R$, on écrit $$P=\prod_{k=1}^{n}(X-\lambda_k)$$ Alors pour tout $z\in\C$ : $$\lvert P(z)\rvert=\prod_{k=1}^{n}\lvert z-\lambda_k\rvert \geqslant \prod_{k=1}^{n}\lvert\im(z-\lambda_k)\rvert=\prod_{k=1}^{n}\lvert\im(z)\rvert=\lvert\im(z)\rvert ^n$$
		\item Sens réciproque : notons $\lambda_k$ les racines \textit{a priori} complexes de $P$ après l'avoir scindé sur $\C$. On a alors : $$\lvert\im(\lambda_k)\rvert^n\leqslant\lvert P(\lambda_k)\rvert=0$$ Puisque $n\geqslant1$, on en déduit que $\im(\lambda_k)=0$ et donc que $\lambda_k\in\R$, si bien que $P$ est scindé sur $\R$.
	\end{itemize}
\end{proof}

\begin{ex}[Moulin.Pol.13]{Des décompositions}
	On pose $P=(X+1)^7-X^7-1\in\R[X]$.
	\begin{enumerate}
		\item Calculer $P(j)$ et en déduire la décomposition en facteurs irréductibles de $P$ dans $\R[X]$.
		\item Donner la décomposition en facteurs irréductibles dans $\R[X]$ de $$\frac{(X^3-1)^4}{P^2}$$
	\end{enumerate}
\end{ex}

\begin{ex}[Châteaux]{Théorème de Gauss-Lucas et applications}
	\begin{enumerate}
		\item Soit $P$ un polynôme complexe non constant. Démontrer que les racines du polynôme $P'$ sont situées dans l'enveloppe convexe des racines de $P$, c'est-à-dire sont des barycentres à coefficients positifs des racines de $P$. Indication : on pourra considérer $P'/P$.
		\item En déduire que si $P$ est scindé dans $\R$, alors $P'$ l'est également, puis que si toute racine de $P$ est de partie réelle positive, alors il en va de même pour $P'$.
	\end{enumerate}
\end{ex}

\begin{ex}[Morlot]{Pas de racine multiple}
	Soit $n \ge 2$ un entier naturel. On considère le polynôme $$P = \sum_{k=0}^{n}\frac{X^k}{k!}$$ Le polynôme $P$ admet-il une racine multiple ?
\end{ex}
\begin{proof}
	La réponse est non. Pour le prouver, raisonner par l'absurde et considérer une racine multiple $z$ de $P$. Par caractérisation de la multiplicité des racines, $z$ est racine de $P'$. Or, on a immédiatement : $$P'=\sum_{k=0}^{n-1}\frac{X^k}{k!}$$ En soustrayant les égalités $P(z)=0$ et $P'(z)=0$, on obtient que $\displaystyle \frac{z^n}{n!}=0$. Par intégrité de $\C$, $z=0$. Or, $z$ n'est clairement pas racine de $P$ car $P(0)=1$. Donc 0 n'est même pas racine mais est racine multiple, ce qui est \textbf{absurde}. Ainsi, pas de racine multiple pour $P$. On remarquera que cet exercice fonctionne car $n \ge 2$ donc $n-1 \ge 1$ donc on a le "droit" à $P'$ quand on applique la caractérisation de la multiplicité des racines.
\end{proof}

\begin{ex}[Moulin.Pol.14]{$\Delta\star$}
	Soit $a\in\R$ et $n\in\N\st$.
	\begin{enumerate}
		\item Résoudre dans $\C$ l'équation $(z+1)^n=e^{2ina}$.
		\item En déduire $\displaystyle\prod_{k=0}^{n-1}\sin\left(a+\frac{k\pi}{n}\right)$ puis $\displaystyle\prod_{k=1}^{n-1}\sin\left(\frac{k\pi}{n}\right)$.
	\end{enumerate}
\end{ex}

\begin{ex}[Moulin.Pol.15]{$\star\star$ Un produit moche}
	Calculer $\displaystyle\prod_{k=1}^{n}\left(5-4\cos\left(\frac{k\pi}{n}\right)\right)$.
\end{ex}
\begin{proof}
	On remarque que $$X^{2n}-1=\prod_{k=1}^{2n}\left(X-\exp\left(i\frac{2k\pi}{2n}\right)\right)=(X^2-1)\prod_{k=1}^{n-1}\left(X-\exp\left(i\frac{2k\pi}{2n}\right)\right)\prod_{k=n+1}^{2n-1}\left(X-\exp\left(i\frac{2k\pi}{2n}\right)\right)$$ Ensuite, dans les produits, on regroupe les termes d'indices : $1$ et $2n-1$, puis $2$ et $2n-2$, etc. pour obtenir : $$X^{2n}-1=(X^2-1)\prod_{k=1}^{n-1}\left(X^2-2\cos\left(\frac{k\pi}{n}\right)X+1\right)$$ On substitue $2$ à $X$ pour obtenir : $$4^n-1=3\prod_{k=1}^{n-1}\left(5-4\cos\left(\frac{k\pi}{n}\right)\right)$$ Le facteur correspondant à $k=n$ valant $9$, on obtient finalement : $$\prod_{k=1}^{n}\left(5-4\cos\left(\frac{k\pi}{n}\right)\right)=3(4^n-1)$$
\end{proof}

\begin{ex}[Moulin.Pol.16]{}
	Déterminer les $P\in\C[X]$ tels que $P(X+1)P(X)=P(X^2)$.
\end{ex}
\begin{proof}
	S'intéresser aux racines potentielles de ce polynôme. Elle doivent être nulles ou de module $1$, puis on peut même prouver que ce doivent des racines de l'unité particulières.
\end{proof}

\begin{ex}[Moulin.Pol.17]{$\star$}
	Soit $P\in\C[X]$ de degré $n\geqslant 1$.
	\begin{enumerate}
		\item Montrer que l'on peut écrire $P=P_1+P_2$ avec les $P_i$ de degré $n$ et de racines complexes de module inférieur à $1$.
		\item $\star\star$ Montrer que l'on peut écrire $P=P_1+P_2+P_3+P_4$ avec les $P_i$ de degré $n$ et de racines complexes de module $1$. Indication : on pourra utiliser l'application $$z\mapsto\frac{z-u}{1-z\overline{u}}$$ lorsque $\lvert u\rvert <1$.
	\end{enumerate}
\end{ex}
\begin{proof}
	Voir la première partie de Maths C 2021.
\end{proof}

\begin{ex}[Moulin.Pol.18]{Un lemme}
	Pour une fraction rationnelle non nulle $F\in\C(X)$, on note $n(F)$ le nombre de ses racines (distinctes), et $p(F)$ le nombre de ses pôles.
	\begin{enumerate}
		\item Soit $P\in\C[X]$ non constant. Montrer que $n(P)+n(P-1)>\deg(P)$.
		\item $\star$ Soit $F\in\C(X)$ non constante. Montrer que $n(F)+n(F-1)+p(F)>a$, où $a$ désigne le degré du numérateur de $F$ dans son écriture sous forme irréductible.
	\end{enumerate}
\end{ex}
\begin{proof}
	Exercice qui sert à démontrer plusieurs résultats.
	\begin{enumerate}
		\item On rappelle que $P_1=P\wedge P'$ est de degré $\deg(P)-n(P)$. De même, $P_2=(P-1)\wedge P'$ est de degré $\deg(P-1)-n(P-1)=\deg(P)-n(P-1)$. Or, $P_1\wedge P_2=1$ et ces deux polynômes divisent $P'$, donc leur produit divise $P'$. Ainsi : $$\deg(P_1)+\deg(P_2)\leqslant\deg(P')=\deg(P)-1$$ On en déduit : $$n(P)+n(P-1)\geqslant \deg(P)+1$$
		\item On écrit $F=P/Q$ sous forme irréductible. Travailler sur $F'$ : les racines de $P$, $Q$ et $P-Q$ sont racines du dénominateur de $F'$. Puis utiliser le même genre d'argument que précédemment.
	\end{enumerate}
\end{proof}

\begin{ex}[Moulin.Pol.19]{Applications de l'exercice précédent}
	\begin{enumerate}
		\item Soit $P$ et $Q$ deux polynômes non constants dans $\C[X]$ tels que $P\inv(\{0\})=Q\inv(\{0\})$ et $P\inv(\{1\})=Q\inv(\{1\})$. Montrer que $P=Q$.
		\item Théorème de Mason : soit $A$, $B$ et $C$ trois polynômes non nuls premiers entre eux dans leur ensemble, non tous constants et tels que $A+B+C=0$. On note $n$ le maximum des degrés de $A$, $B$ et $C$. Montrer que $ABC$ possède au moins $n+1$ racines distinctes.
		\item Soit $m\geqslant 3$ entier et $(A,B,C)\in\C[X]^3$ un triplet de polynômes premiers entre eux dans leur ensemble tel que $A^m+B^m=C^m$. Montrer que $A$, $B$ et $C$ sont tous constants.
	\end{enumerate}
\end{ex}
\begin{proof}
	Plein de jolies questions qui sont à elles seules des exercices d'oraux !
	\begin{enumerate}
		\item $P-Q=(P-1)-(Q-1)$ possède au moins $n(P)+n(P-1)$ racines distinctes, soit strictement plus de racines que $\deg(P)$ d'après l'exercice précédent. Or, $\deg(P-Q)\leqslant \deg(P)$ donc $P-Q=0$ et $P=Q$. \\
		Remarque : l'énoncé fonctionne pour $a$ et $b$ deux complexes distincts.
		\item $A$, $B$ et $C$ sont deux à deux premiers entre eux. En réécrivant la relation sous la forme : $$1+\frac{B}{A}=\frac{-C}{A}$$ et en utilisant la question $2$ de l'exercice précédent avec $F=1+B/A$, on trouve : $n(C)+n(B)+n(A)>\deg(C)$. Par symétrie de $A$, $B$ et $C$, cela est encore vrai avec $\deg(A)$ et $\deg(B)$, si bien que cela est vrai pour le maximum des degrés. Observer pour conclure que le nombre de racines distinctes de $ABC$ est $n(A)+n(B)+n(C)$ car $A$, $B$ et $C$ sont deux à deux premiers entre eux.\\
		Alternativement, on peut remarquer que $AB'-BA'=BC'-CB'=CA'-AC'=P$ est un polynôme constant. Or, $A$, $B$ et $C$ sont deux à deux premiers entre eux (en prenant un facteur commun irréductible de deux polynômes, il divise le troisième donc est constant par hypothèse). Comme au moins $A'$, $B'$ ou $C'$ est non nul, $P$ est non nul sinon $A$ diviserait $A'$ par exemple (en supposant $A$ non constant), ce qui est absurde. En notant $Q_1$ le PGCD de $Q$ et $Q'$ pour $Q=A$ puis $B$ puis $C$, $A_1$, $B_1$ et $C_1$ sont premiers entre eux et divisent $P$ donc la somme de leurs degrés est inférieure au degré de $P$. On utilise ensuite le fait que le degré de $A_1$ est supérieur au degré de $A$ moins le nombre de racines distinctes de $A$. De plus, $\deg(P)\leqslant\deg(A)+\deg(B)-1$. En substituant, en trouve que $\deg(C)$ est inférieur ou égal à la somme du nombre de racines distinctes de $A$, $B$ et $C$ moins $1$, et la somme de ces nombres de racines distinctes est exactement égale au nombre de racines distinctes de $ABC$ car $A$, $B$ et $C$ sont deux à deux premiers entre eux.
		\item De tels polynômes sont alors deux à deux premiers entre eux : en effet, si $P$ irréductible divise $A$ et $B$, alors il divise $C^m$ donc il divise $C$, donc il est constant. De même, on obtient le fait que ces polynômes sont deux à deux premiers entre eux. Par l'absurde, on suppose que $A$, $B$ ou $C$ n'est pas constant et on note $n$ le degré maximal (supérieur à $1$) de l'un de ces $3$ polynômes. Alors, d'après la question $2$, $-(ABC)^m$ possède au moins $nm+1$ racines distinctes. Donc $ABC$ possède $nm+1$ racines distinctes. Donc $A$, $B$ ou $C$ est nul. Donc $A$, $B$ et $C$ ne sont plus premiers entre eux dans leur ensemble. C'est absurde, donc $A$, $B$ et $C$ sont tous constants.
	\end{enumerate}
\end{proof}

\begin{ex}[Moulin.Pol.20]{Autour de l'image de $\U$}
	\begin{enumerate}
		\item Déterminer tous les polynômes $P\in\C[X]$ tels que $\P(\U)\subset\U$ puis tels que $\P(\U)=\U$.
		\item $\star$ Déterminer toutes les fractions rationnelles $F\in\C(X)$ n'ayant pas de pôle dans $\U$ et telles que $F(\U)\subset\U$.
	\end{enumerate}
\end{ex}

\begin{ex}[Moulin.Pol.21]{Contrôle du module d'une racine particulière}
	Soit $p_1,\ldots,p_n$ des réels positifs non tous nuls.
	\begin{enumerate}
		\item Montrer que $P=X^n-p_1X^{n-1}-\cdots-p_n$ admet une unique racine dans $\R_+\st$.
		\item Soit $a_1,\ldots,a_n$ des nombres complexes non tous nuls, $z_0$ une racine de $X^n+a_1X^{n-1}+\cdots+a_n$ et $e$ la racine strictement positive du polynôme $P$ où on a pris $p_i=\lvert a_i\rvert$ et $P$ est défini dans la première question. Montrer que $\lvert z_0\rvert\leqslant e$.
		\item Soit $b_1,\ldots,b_n$ des réels strictement positifs de somme inférieure ou égale à $1$. Montrer que $$\lvert z_0\rvert\leqslant\max_{i\in\llbracket1,n\rrbracket}\left\lvert\frac{a_i}{b_i}\right\rvert^{1/i}$$
	\end{enumerate}
\end{ex}

\begin{ex}[Moulin.Pol.22]{Un polynôme scindé à racines simples}
	Soit $n\in\N\st$ et $P(X)=nX^n-\displaystyle\sum_{k=0}^{n-1}X^k\in\C[X]$.
	\begin{enumerate}
		\item Soit $z$ une racine de $P$ distincte de $1$. Montrer que $\lvert z\rvert <1$.
		\item Montrer que les racines de $P$ sont simples.
	\end{enumerate}
\end{ex}
\begin{proof}
	Remarquer que $1$ est racine. On peut ensuite écrire $P=(X-1)Q'$ où $Q=1+X+\cdots+X^n$ et étudier $Q'$ en observant que $Q=(X^{n+1}-1)/(X-1)$.
\end{proof}

\begin{ex}[Moulin.Pol.23]{$\star$ Théorème d'Eneström-Kakeya}
	Soit $(a_0,\ldots,a_n)\in\R^{n+1}$ et $P=\displaystyle\sum_{k=0}^{n}a_kX^k$.
	\begin{enumerate}
		\item Montrer que si $a_0\geqslant a_1\geqslant\cdots\geqslant a_n>0$, alors toutes les racines complexes de $P$ sont de module au moins $1$. Indication : on pourra poser $\alpha_k=a_k-a_{k+1}$.
		\item En déduire que si l'on suppose seulement les $a_k$ strictement positifs, alors toute racine $z$ de $P$ vérifie : $$\min_{0\leqslant k<n}\frac{a_k}{a_{k+1}}\leqslant \lvert z\rvert\leqslant\max_{0\leqslant k<n}\frac{a_k}{a_{k+1}}$$
	\end{enumerate}
\end{ex}
\begin{proof}
	Exercice Cassini...
	\begin{enumerate}
		\item On montre à l'aide de l'inégalité triangulaire : $$\forall z\in\C,\ \lvert (1-z)P(z)\rvert\geqslant a_0-(\alpha_0\lvert z\rvert+\cdots+\alpha_{n-1}\lvert z\rvert^{n-1})$$ Donc si $\lvert z\rvert >1$, alors $\lvert (1-z)P(z)\rvert>0$.
		\item On note $m$ et $M$ les minimum et maximum mentionnés dans l'énoncé. Appliquer la question $1$ aux polynômes $P(mX)$ et $X^nP(M/X)$ pour conclure.
	\end{enumerate}
\end{proof}

\begin{ex}[Moulin.Pol.24]{$\star\star$}
	Soit $(a_0,\ldots,a_n)\in\R^{n+1}$ et $f:\theta\mapsto\displaystyle\sum_{k=0}^{n}a_k\cos(k\theta)$. On suppose $\forall\theta\in\R,\ f(\theta)\geqslant0$. Montrer qu'il existe $(b_0,\ldots,b_n)\in\C^{n+1}$ tel que $$\forall\theta\in\R,\ f(\theta)=\left\lvert\sum_{k=0}^{n}b_ke^{ik\theta}\right\rvert^2$$
\end{ex}
\begin{proof}
	Exercice redoutable. De façon plus générale, on peut prendre : $$f:\theta\mapsto\sum_{k=-n}^{n}c_ke^{ik\theta}$$ avec les $c_k$ complexes et on suppose $c_n$ non nul (sinon on réduit simplement le problème). On écrit $f(\theta)=e^{-in\theta}P(e^{i\theta})$ où $P\in\C_{2n}[X]$ avec $P(0)\ne 0$. Puisque $f$ est à valeurs réelles, on observe que nécessairement, on a : $$c_{-k}=\overline{c_k}$$ Cela permet de montrer, pour tout $z\in\C\st$ : $$P(z)=z^{2n}\displaystyle\overline{P\left(\frac{1}{\overline{z}}\right)}$$ Si $e^{ix}$ est une racine de $P$ de module $1$ et de multiplicité $m$, on montre en utilisant un arc-moitié que l'on peut écrire : $$\forall\theta\in\R,\ f(\theta)=\sin^m\left(\frac{\theta-x}{2}\right)g(\theta)$$ où $g(x)\ne 0$. Comme $f$ est à valeurs dans $\R_+$, on en déduit même que $m$ est un entier pair. Ensuite, puisque $0$ n'est pas racine de $P$, on montrer que pour toute racine $z$ de module différent de $1$, $1/\overline{z}$ est aussi une racine de $P$ de même multiplicité. On scinde ensuite $P$ sur $\C$. On établit alors l'égalité : $$\forall(\theta,z)\in\R\times\C\st,\ \lvert (e^{i\theta}-z)(e^{i\theta}-1/\overline{z})\rvert=\frac{\lvert e^{i\theta}-z\rvert^2}{\lvert z\rvert}$$ On remarque alors que $f(\theta)=\lvert f(\theta)\rvert=\lvert P(e^{i\theta})\rvert$ et on assemble ce qui précède pour conclure.
\end{proof}

\begin{ex}[Moulin.Pol.25]{$\star\star$}
	Soit $P=\displaystyle\sum_{k=0}^{n}a_kX^k\in\C[X]$ dont toutes les racines ont une partie imaginaire strictement positive. Montrer que $Q=\displaystyle\sum_{k=0}^{n}\text{Re}(a_k)X^k$ est scindé sur $\R$. Indication : comparer $\lvert P(z)\rvert$ et $\lvert P(\overline{z})\rvert$ pour $z$ non réel.
\end{ex}
\begin{proof}
	Tout d'abord, on suit l'indication : après avoir scindé $P$ sur $\C$, on se ramène à l'étude de $\lvert z-\lambda\rvert$ et $\lvert \overline{z}-\lambda\rvert$ où $\text{Im}(\lambda)>0$. Sans perte de généralité, on suppose $\text{Im}(z)>0$. Par un calcul utilisant les parties réelles et imaginaires, on montre aisément $$\lvert z-\lambda\rvert<\lvert \overline{z}-\lambda\rvert$$ Avoir une intuition géométrique de l'inégalité précédente peut aider. On en déduit que si $z$ est de partie imaginaire strictement positive, alors : $$\lvert P(z)\rvert <\lvert P(\overline{z})\rvert$$ De façon générale, pour $z$ complexe non réel, on a $\lvert P(z)\rvert \ne \lvert P(\overline{z})\rvert$. \\
	Ensuite, on note $\overline{P}$ le polynôme dont les coefficients sont les conjugués des coefficients de $P$. Alors, d'après les formules d'Euler, on a : $$Q=\frac{P+\overline{P}}{2}$$ Mais $Q$ est scindé sur $\C$. Soit $z$ une racine complexe de $Q$. Alors, $P(z)=-\overline{P}(z)$. Donc $$\lvert P(z)\rvert=\lvert -\overline{P}(z)\rvert=\lvert\overline{P}(z)\rvert=\lvert\overline{\overline{P}(z)}\rvert=\lvert P(\overline{z})\rvert$$ D'après l'étude précédente, par contraposée, $z$ est nécessairement réel. Donc $Q$ est scindé sur $\R$.
\end{proof}

\begin{ex}[Moulin.Pol.26]{}
	Soit $P=a_0+\cdots+a_nX^n\in\Z[X]$ avec $n\in\N\st$ et $a_n\ne 0$. On pose $M=\displaystyle\sup_{i\in\llbracket0,n-1\rrbracket}\lvert a_i/a_{n}\rvert$.
	\begin{enumerate}
		\item Montrer que toute racine de $P$ est de module strictement inférieur à $M+1$.
		\item On suppose qu'il existe un entier naturel $k\geqslant M+2$ tel que $P(k)$ soit premier. Montrer que $P$ est irréductible dans $\Z[X]$.
	\end{enumerate}
\end{ex}
\begin{proof}
	Cet exercice donne un critère d'irréductibilité (s'il en fallait encore..) des polynômes à coefficients entiers.
	\begin{enumerate}
		\item Soit $z$ une racine de $P$. Si $\lvert z\rvert<1$, le résultat est immédiat. Si un des $a_i$ pour $i\in\llbracket0,n-1\rrbracket$ est non nul, le résultat est immédiat pour $\lvert z\rvert=1$. Si tous les $a_i$ sont nuls pour $i\in\llbracket0,n-1\rrbracket$, les seules racines de $P$ sont nulles, donc le résultat vaut en toute généralité pour $\lvert z\rvert\leqslant 1$. On suppose désormais $\lvert z\rvert >1$. Alors : $$-a_nz^n=a_0+\cdots+a_{n-1}z^{n-1}$$ donc après division par $a_n$, utilisation de l'inégalité triangulaire et majoration des quotients par $M$, on obtient : $$\lvert z\rvert^n\leqslant M(1+\cdots+\lvert z\rvert^{n-1})=M\frac{\lvert z\rvert^n-1}{\lvert z\rvert-1}\leqslant M\frac{\lvert z\rvert^n}{\lvert z\rvert-1}$$ En simplifiant par $\lvert z\rvert^n$ et en isolant $\lvert z\rvert$, on obtient le résultat.
		\item Soit $A$ et $B$ dans $\Z[X]$ tels que $P=AB$. $P(k)=A(k)B(k)$ est premier et $A(k)$ et $B(k)$ sont entiers, donc $A(k)=\pm1$ ou $B(k)=\pm1$. Sans perte de généralité, on suppose $A(k)=\pm1$, si bien que le coefficient dominant de $A$ vaut $\pm1$. Mais alors, en scindant $A$ sur $\C$ si $A$ est non constant, toutes les racines de $A$ sont racines de $P$, donc de module strictement inférieur à $M+1$ d'après la question $1$. Donc tous les facteurs de $A(k)$ sont module supérieur à $k-(M+1)>1$ par l'inégalité triangulaire inversée. Mais cela est impossible, car $A(k)=\pm1$ est de module $1$. Donc $A$ est constant, et $P$ est irréductible dans $\Z[X]$.
	\end{enumerate}
\end{proof}

\section{Racines réelles}

\begin{ex}[Moulin.Pol.27]{$\Delta$}
	Soit $P \in \R[X]$ un polynôme scindé.
	\begin{enumerate}
		\item Supposons que $P$ est à racines simples. Montrer que $P'$ est scindé.
		\item Le résultat est-il toujours valable lorsque les racines de $P$ ne sont plus nécessairement simples ?
	\end{enumerate}
\end{ex}
\begin{proof}
	Le maître mot : théorème de Rolle. 
	\begin{enumerate}
		\item Ordonner les racines de $P$ et utiliser le théorème de Rolle entre deux racines successives de $P$.
		\item Réutiliser l'idée précédente en faisant attention à la multiplicité potentielle des racines. Le résultat reste valable.
	\end{enumerate}
\end{proof}

\begin{ex}[Moulin.Pol.28]{Comportement des racines de $P'$}
	\begin{enumerate}
		\item Montrer que pour tout polynôme $P$ réel de degré $n$ à racines $a_1<\ldots<a_n$, le polynôme $P'$ a $n-1$ racines vérifiant $a_1<b_1<a_2<\ldots<a_{n-1}<b_{n-1}<a_n$.
		\item Montrer que pour tout $i\in\llbracket1,n\rrbracket$ : $$n\prod_{j=1}^{n-1}(a_i-b_j)=\prod_{j\ne i}(a_i-a_j)$$
		\item En déduire : $$\forall i\in\llbracket1,n-1\rrbracket,\ a_i+\frac{a_{i+1}-a_i}{n}\leqslant b_i\leqslant a_{i+1}-\frac{a_{i+1}-a_i}{n}$$
	\end{enumerate}
\end{ex}
\begin{proof}
	Faire l'exercice précédent avant.
	\begin{enumerate}
		\item Voir la première question de l'exercice précédent.
		\item Évaluer $P'$ en $a_i$ en écrivant $P'$ de deux manières différentes : en le décomposant en produit d'irréductibles pour le premier membre, et en dérivant la décomposition en irréductibles de $P$ pour le membre de droite. Pour terminer, simplifier par le coefficient dominant de $P$.
		\item Utiliser l'égalité précédente et basculer tous les facteurs du produit de gauche à droite sauf $a_i-b_i$ (resp. $a_{i+1}-b_i$) Utiliser ensuite des majorations ou des minorations triviales comme : $$\frac{a_1-a_2}{a_1-b_1}\geqslant 1$$ pour majorer ou minorer grossièrement $n-2$ facteurs : il ne reste alors plus que le $n$ et le facteur qui nous intéresse.
	\end{enumerate}
\end{proof}

\begin{ex}[Moulin.Pol.29]{$\Delta$}
	Montrer que $L_n=\left((X^2-1)^n\right)^{(n)}$ est scindé à racines simples sur $\R$.
\end{ex}
\begin{proof}
	Itérer l'utilisation du théorème de Rolle.
\end{proof}

\begin{ex}[Moulin.Pol.30]{Une suite de polynômes scindés}
	Montrer qu'il existe une suite $(a_n)\in(\R\st)^\N$ telle que pour tout $n\in\N\st$, $P_n=\displaystyle\sum_{k=0}^{n}a_kX^k$ soit scindé à racines simples sur $\R$.
\end{ex}
\begin{proof}
	Construire la suite par récurrence en prenant les coefficients suffisamment petits pour scinder encore sans trop modifier les racines.
\end{proof}

\begin{ex}[Moulin.Pol.31]{$\star$}
	Soit $P\in\R[X]$ admettant exactement $p>0$ coefficients non nuls.
	\begin{enumerate}
		\item Montrer que $P$ admet au plus $2p-1$ racines réelles distinctes.
		\item Montrer que cette majoration est optimale pour tout $p$.
	\end{enumerate}
\end{ex}

\begin{ex}[Moulin.Pol.32]{$\star$}
	Soit $P\in\R[X]$ scindé sur $\R$.
	\begin{enumerate}
		\item Montrer que pour tout $\alpha\in\R$, le polynôme $P'-\alpha P$ est scindé sur $\R$.
		\item Montrer que si $Q=\displaystyle\sum_{k=0}^{n}b_kX^k$ est aussi scindé sur $\R$, alors il en est de même pour $\displaystyle\sum_{k=0}^{n}b_kP^{(k)}$.
	\end{enumerate}
\end{ex}
\begin{proof}
	La deuxième question se déduit assez facilement de la première.
	\begin{enumerate}
		\item Soit $\alpha\in\R$. On considère $f:t\mapsto e^{-\alpha t}P(t)$ définie sur $\R$. $f$ est dérivable sur $\R$ et on a : $$\forall t\in\R,\ f'(t)=e^{-\alpha t}(P'(t)-\alpha P(t))$$ Donc $f'(t)=0\iff P'(t)-\alpha P(t)=0$. Or, si $x$ est une racine de $m\geqslant 1$ de $P$, alors $x$ est une racine de multiplicité $m-1$ de $P'$. Entre deux racines consécutives de $P$ (donc de $f$), on trouve un zéro de $f'$ avec le théorème de Rolle. Il en manque alors une. On la trouve aussi avec une généralisation du théorème de Rolle entre une racine extrémale de $P$ et $\pm\infty$ (en fonction de là où $f$ tend vers $0$). Ou alors, pour la dernière, on remarque qu'on a déjà $n-1$ et qu'on est de degré $n$ (et si $\alpha=0$ c'est trivial) donc il y a nécessairement une dernière racine.
		\item On procède par récurrence sur le degré de $Q$ en utilisant la première question et en regardant $Q(D)$ où $D$ est l'endomorphisme de dérivation de $\R[X]$.
	\end{enumerate}
\end{proof}

\section{Polynômes d'interpolation}

\begin{ex}[Morlot]{$\circ$ Coefficients rationnels}
	Soit $P \in \C[X]$ tel que $\forall x \in \Q, \ P(x)\in\Q$. Montrer que $P$ est à coefficients rationnels.
\end{ex}
\begin{proof}
	Considérer $n$ le degré du polynôme puis utiliser les polynômes interpolateurs de Lagrange. associés à $0,1,\ldots,n$.	
\end{proof}

\begin{ex}[Moulin.Pol.33]{$\star$}
	Soit $d$ et $n$ deux entiers tels que $d\geqslant n\geqslant 1$, et $F\in\C_n[X]$. On suppose qu'il existe un sous-ensemble $A$ de $\U_{d+1}$ de cardinal $n+1$ tel que $\forall z\in A,\ \lvert F(z)\rvert\leqslant 1$. Montrer que $$\max_{z\in\U}\lvert F(z)\rvert\leqslant 2^d$$
\end{ex}

\begin{ex}[Moulin.Pol.34]{}
	Soit $n$ et $k$ deux entiers naturels. On considère des réels $x_1,\ldots,x_n$ deux à deux distincts et des réels $y_1,\ldots,y_n$ quelconques. Pour $P\in\R_{k-1}[X]$, on note $E_P=\{i\in\llbracket1,n\rrbracket\ : \ P(x_i)=y_i\}$.
	\begin{enumerate}
		\item Donner une condition nécessaire et suffisante sur $k$ et $n$ pour que, quel que soit $(y_1,\ldots,y_n)$, il existe $P\in\R_{k-1}[X]$ tel que $\card(E_P)=n$. Y a-t-il alors unicité de $P$ ?
		\item On suppose $k< n$.
		\begin{enumerate}[label=\alph*.]
			\item Montrer qu'il existe au plus un polynôme $P\in\R_{k-1}[X]$ tel que $\displaystyle\card(E_P)\geqslant\frac{n+k}{2}$.
			\item Montrer qu'il existe deux polynômes $Q$ et $R$, avec $R\ne 0$, tels que $$\deg(Q)<\frac{n+k}{2},\ \deg(R)\leqslant\frac{n-k}{2}\ \text{et}\ \forall i\in\llbracket1,n\rrbracket,\ Q(x_i)-y_iR(x_i)=0$$
			\item Montrer que s'il existe $P\in\R_{k-1}[X]$ tel que $\displaystyle\card(E_P)\geqslant\frac{n+k}{2}$, alors $P=Q/R$.
		\end{enumerate}
	\end{enumerate}
\end{ex}

\begin{ex}[Moulin.Pol.35]{$\star$}
	Soit $P=\displaystyle\sum_{j=0}^{n-1}a_jX^j\in\Z[X]$. On suppose que $a_0=1$ et $P(\U_n)\subset\R_+$ où $\U_n$ désigne le groupe des racines $n$-èmes de l'unité. Montrer que les $a_j$ sont dans $\{-1,0,1\}$.
\end{ex}
\begin{proof}
	On écrit, avec les notations habituelles : $$P=\sum_{j=0}^{n-1}P(\omega_j)L_j$$ Or, $L_j=\displaystyle\prod_{\substack{0\leqslant k\leqslant n-1\\ k\ne j}}\frac{X-\omega_k}{\omega_j-\omega_k}$. On va chercher à simplifier l'expression de $L_j$ : pour cela, on utilise le polynôme $X^n-1$ : $$X^n-1=\prod_{k=0}^{n-1}(X-\omega_k)$$ et en dérivant : $$nX^{n-1}=\sum_{k=0}^{n-1}\prod_{l\ne k}(X-\omega_l)$$ Le numérateur de $L_j$ vaut alors : $$\frac{X^n-1}{X-\omega_j}$$ et son dénominateur $n/\omega_j$. En utilisant une factorisation de $a^-b^n$, on trouve alors : $$L_j=\frac{1}{n}\sum_{k=0}^{n}(\omega_j\inv X)^k$$ En sommant, on en déduit alors : $$a_j=\frac{1}{n}\sum_{k=0}^{n-1}P(\omega_k)\omega_k^{-j}$$ Par inégalité triangulaire, comme les $P(\omega_k)$ sont des réels positifs, on obtient : $$\lvert a_j\rvert\leqslant\sum_{k=0}^{n-1}P(\omega_k)$$ Mais en prenant l'égalité précédente pour $j=0$, on trouve $$\sum_{k=0}^{n-1}P(\omega_k)=na_0=n$$ Ainsi : $$\forall j\in\llbracket0,n-1\rrbracket,\ \lvert a_j\rvert\leqslant 1$$ Comme les $a_j$ sont entiers, on en déduit bien : $$\forall j\in\llbracket0,n-1\rrbracket,\ a_j\in\left\{-1,0,1\right\}$$
\end{proof}

\section{Fractions rationnelles}

\begin{ex}[Moulin.Pol.36]{Automorphismes de $\K(X)$}
	\begin{enumerate}
		\item Soit $\Phi$ un endomorphisme de $\K(X)$. Montrer que $\Phi(X)$ est non nul et caractérise entièrement $\Phi$.
		\item Soit $(a,b,c,d)\in\K^4$ tel que $ad-bc\ne 0$. Montrer que $\Phi(X)=\displaystyle\frac{aX+b}{cX+d}$ définit un automorphisme de $\K(X)$.
		\item $\star$ Montrer que tout automorphisme de $\K(X)$ est de cette forme.
	\end{enumerate}
\end{ex}

\begin{ex}[Moulin.Pol.37]{Des sommes}
	Soit $P(X)=\displaystyle\prod_{k=1}^{n}(X-x_k)$, où les $x_k$ sont deux à deux distincts dans $\K$. Calculer : 
	\begin{enumerate}
		\item $\displaystyle\sum_{k=1}^{n}\frac{1}{P'(x_k)}$ ;
		\item $\displaystyle\sum_{k=1}^{n}\frac{x_k}{P'(x_k)}$ ;
		\item $\displaystyle\sum_{k=1}^{n}\frac{1}{x_kP'(x_k)}$ (en supposant que les $x_k$ sont non nuls pour cette dernière).
	\end{enumerate}
\end{ex}

\begin{ex}[Moulin.Pol.38]{$\star$ Inégalité de Bernstein}
	Soit $n\geqslant 1$ un entier. On note $z_1,\ldots,z_n$ les racines de $X^n+1$.
	\begin{enumerate}
		\item Montrer que pour tout $P\in\C_n[X]$ $$XP'=\frac{n}{2}P+\frac{2}{n}\sum_{k=1}^{n}\frac{z_kP(z_kX)}{(z_k-1)^2}$$
		\item Pour $Q\in\C[X]$, on pose $\lVert Q\rVert=\displaystyle\sup_{z\in\U}\lvert Q(z)\rvert$. Montrer que pour tout $P\in\C_n[X]$, $$\lVert P'\rVert\leqslant n\lVert P\rVert$$
	\end{enumerate}
\end{ex}
\begin{proof}
	Pour la première question, remarquer une linéarité en $P$ pour se ramener à la base canonique de $\C_n[X]$. On utilisera, après démonstration, la décomposition en éléments simples de $X^k/(X^n+1)$ pour $k\in\llbracket 0,n\rrbracket$.
\end{proof}

\begin{ex}[Moulin.Pol.39]{Un calcul de somme}
	\begin{enumerate}
		\item Soit $x_0,\ldots,x_n$ des éléments de $\K$ distincts deux à deux. Calculer $$\sum_{k=0}^{n}\frac{x_k^p}{\prod_{i\ne k}(x_k-x_i)}$$ en fonction de $p\in\llbracket0,n\rrbracket$.
		\item $\star$ En déduire la valeur de $\displaystyle\sum_{k=0}^{n}(-1)^kk^p\binom{n}{k}$ pour $p\in\llbracket0,n\rrbracket$.
	\end{enumerate}
\end{ex}
\begin{proof}
	La deuxième question n'est pas si difficile que ça, je ne sais pas trop pourquoi il y a une étoile.
	\begin{enumerate}
		\item On note $P=\displaystyle\prod_{k=0}^{n}(X-x_k)$ et on considère la fraction rationnelle : $$F_p=\frac{X^p}{P}$$ $F_p$ est de degré $p-(n+1)<0$ donc on peut écrire sa décomposition en éléments simples : $$F_p=\sum_{k=0}^{n}\frac{\lambda_{k,p}}{X-x_k}$$ De plus, on sait que : $$\lambda_{k,p}=\frac{x_k^p}{P'(x_k)}=\frac{x_k^p}{\prod_{i\ne k}(x_k-x_i)}$$ Donc : $$F_p=\sum_{k=0}^{n}\frac{x_k^p}{\prod_{i\ne k}(x_k-x_i)}\times\frac{1}{X-x_k}$$ On multiplie par $X$, puis on substitue $x$ qui tend vers $+\infty$ pour obtenir la valeur de la somme. Alors : 
		\begin{itemize}
			\item si $p\in\llbracket0,n-1\rrbracket$, la somme vaut $0$ car $xF_p(x)$ tend vers $0$ quand $x$ tend vers $+\infty$ ;
			\item si $p=n$, la somme vaut $1$ car $xF_p(x)$ tend vers $1$ quand $x$ tend vers $+\infty$.
		\end{itemize}
		\item Prendre $x_k=k$ et utiliser le résultat de la question précédente après simplification des termes. La somme vaut alors :
		\begin{itemize}
			\item $0$ si $p\in\llbracket0,n-1\rrbracket$ ;
			\item $(-1)^nn!$ si $p=n$.
		\end{itemize}
	\end{enumerate}
\end{proof}

\begin{ex}[Moulin.Pol.40]{}
	Soit $P=\displaystyle\sum_{k=0}^{n}a_kX^k\in\R[X]$ scindé.
	\begin{enumerate}
		\item Montrer : $\forall x\in\R,\ P(x)P''(x)\leqslant P'(x)^2$.
		\item En déduire : $\forall k\in\llbracket1,n-1\rrbracket,\ a_{k-1}a_{k+1}\leqslant a_k^2$.
		\item $\star$ Montrer plus précisément : $\forall x\in\R,\ nP(x)P''(x)\leqslant(n-1)P'(x)^2$.
	\end{enumerate}
\end{ex}
\begin{proof}
	On utilise ici principalement la forme $P'/P$ pour $P$ scindé. On note $x_1,\ldots,x_n$ les racines de $P$ comptées avec multiplicité.
	\begin{enumerate}
		\item Si $x$ est une racine de $P$, le résultat est immédiat. Sinon : $$\frac{P'(x)}{P(x)}=\sum_{k=1}^{n}\frac{1}{x-x_k}$$ En dérivant cette expression sur les intervalles ouverts ne contenant pas de racine de $P$, on trouve que pour $x$ non racine de $P$ : $$\frac{P''(x)P(x)-P'(x)^2}{P(x)^2}=-\sum_{k=1}^{n}\frac{1}{(x-x_k)^2}\leqslant 0$$
		\item Prendre $x=0$ dans l'inégalité précédente donne : $a_0\times 2a_2\leqslant a_1^2$ puis on divise par $2$ et on utilise le fait que $1/2\leqslant 1$. Ensuite, appliquer ce même résultat aux coefficients de $P^{(k-1)}$ qui est aussi scindé (cf. un autre exercice).
		\item Il s'agit ici de raffiner l'inégalité obtenue à la question $1$. Pour cela, on observe que d'après l'inégalité de Cauchy-Schwarz, pour $x$ non racine de $P$ : $$\left(\frac{P'(x)}{P(x)}\right)^2=\left(\sum_{k=1}^{n}\frac{1}{x-x_k}\times 1\right)^2\leqslant n\sum_{k=1}^{n}\frac{1}{(x-x_k)^2}$$ Dès lors, pour $x$ non racine de $P$ : $$\frac{P''(x)P(x)-P'(x)^2}{P(x)^2}\leqslant-\frac{1}{n}\left(\frac{P'(x)}{P(x)}\right)^2$$ En multipliant par $nP'(x)^2$ en basculant un $P'(x)^2$ de l'autre côté de l'inégalité, on obtient l'inégalité souhaitée.
	\end{enumerate}
\end{proof}

\begin{ex}[Moulin.Pol.41]{$\star$ Théorème de Jensen}
	Soit $P\in\R[X]$ de degré supérieur ou égal à $2$. On note $x_1,\ldots,x_p$ les racines réelles de $P$ et $z_1,\overline{z_1},\ldots,z_q,\overline{z_q}$ les racines de $P$ dans $\C\setminus\R$. Si $k\in\llbracket 1,q\rrbracket$, on note $D_k$ le disque de diamètre $[z_k,\overline{z_k}]$. Montrer que les racines non réelles de $P'$ sont dans $D_1\cup\ldots\cup D_q$.
\end{ex}

\begin{ex}[Moulin.Pol.42]{$\star$}
	Soient $n\in\N\st$, $P\in\C[X]$ de degré $n$ et $Q=2XP'-nP$.
	\begin{enumerate}
		\item On suppose que les racines de $P$ sont toutes dans $\U$. Montrer qu'il en est de même pour celles de $Q$.
		\item Soit $R\in\R_+$. On suppose que toutes les racines de $P$ sont de module $R$. Que dire de celles de $Q$ ?
	\end{enumerate}
\end{ex}
\begin{proof}
	Exercice donné en colle par FMC à l'époque.
	\begin{enumerate}
		\item C'est la question difficile. La deuxième s'en déduit aisément. L'astuce réside essentiellement dans le fait d'écrire $XP'/P$ de deux façons différentes. On note $\lambda_k$ les racines de $P$ comptées avec multiplicité. Alors, d'après le cours : $$\frac{P'}{P}=\sum_{k=1}^{n}\frac{1}{X-\lambda_k}$$ Si bien que : $$\frac{XP'}{P}=\sum_{k=1}^{n}\frac{X}{X-\lambda_k}=n+\sum_{k=1}^{n}\frac{\lambda_k}{X-\lambda_k}$$ Ensuite, on considère une $z$ racine de $Q$ non racine de $P$ (sinon le résultat est trivial). Comme : $$\frac{Q}{2P}=\frac{XP'}{P}-\frac{n}{2}$$ on en déduit : $$\frac{n}{2}=\begin{cases}
			\displaystyle\sum_{k=1}^{n}\frac{z}{z-\lambda_k}=\sum_{k=1}^{n}\frac{\lvert z\rvert^2-z\overline{\lambda_k}}{\lvert z-\lambda_k\rvert^2}\\
			\displaystyle n+\sum_{k=1}^{n}\frac{\lambda_k}{z-\lambda_k}=n+\sum_{k=1}^{n}\frac{\overline{z}\lambda_k-1}{\lvert z-\lambda_k\rvert^2}
		\end{cases}$$ En passant aux conjugués et en utilisant la deuxième ligne d'égalités précédente : $$-\frac{n}{2}=\sum_{k=1}^{n}\frac{z\overline{\lambda_k}-1}{\lvert z-\lambda_k\rvert^2}$$ En sommant cette égalité avec la première égalité dans dans les deux précédentes, on trouve : $$0=(\lvert z\rvert^2-1)\underbrace{\sum_{k=1}^{n}\frac{1}{\lvert z-\lambda_k\rvert^2}}_{>0}$$ donc $\lvert z\rvert^2-1=0$ et on en déduit $\lvert z\rvert=1$.
		\item Si $r=0$, alors $P=\lambda X^n$ pour $n\in\N\st$ et $\lambda\ne 0$, et alors $Q=n\lambda X^n$ a toutes ses racines nulles donc de module $0$. \\
		Si $r>0$, on pose $\tilde{P}=P(X/r)$ qui a toutes ses racines de module $1$, et on se ramène au résultat de la première question.
	\end{enumerate}
\end{proof}

\section{Arithmétique des polynômes}

\begin{ex}{Restes de divisons euclidiennes}
	Soit $P \in \K[X]$ et $(a,b) \in\K^2$ distincts. 
	\begin{enumerate}
		\item Déterminer le reste de la division euclidienne de $P$ par $(X-a)(X-b)$
		\item Déterminer le reste de la division euclidienne de $P$ par $(X-a)^2$
	\end{enumerate}
\end{ex}
\begin{proof}
	Quel est le degré du reste ? Utiliser des évaluations précises et/ou la dérivation suivie d'une évaluation. Ne pas tenter l'algorithme sur un polynôme quelconque dans un corps quelconque !!! On va plutôt utiliser le fait que le reste est à chaque fois de la forme $\lambda X + \mu$ avec $(\lambda,\mu) \in \K^2$.
	\begin{enumerate}
		\item On peut fixer $Q \in\K[X]$ tel que $P=(X-a)(X-b)Q + (\lambda X + \mu)$. En évaluant en $a$ et en $b$, on obtient $\lambda a + \mu =P(a)$ et $\lambda b + \mu = P(b)$. De simples combinaisons linéaires permettent de conclure.
		\item Là encore, on peut fixer $Q \in\K[X]$ tel que $P=(X-a)^2Q + (\lambda X + \mu)$. Toutefois, on ne peut plus réaliser deux évaluations. Cependant, une évaluation en $a$ fournit déjà
	\end{enumerate}
\end{proof}

\begin{ex}{$\star\star$ Caractérisation d'un corps}
	Soit $A$ un anneau. Montrer que $A$ est un corps si, et seulement si, $A[X]$ est principal.
\end{ex}
\begin{proof}
	Les sens direct est un théorème de cours. Réciproquement, supposons que $A[X]$ soit principal. Il ne nous reste plus qu'à montrer l'inversibilité des éléments non nuls. Soit $a\in A\setminus\{0\}$. On considère $I$ l'idéal engendré par $a$ et $X$. Alors il existe $P\ne 0$ tel que $I=P\times A[X]$. Puisque $a\in I$, $P\mid a$ donc $P$ constant et, abusivement, $P=p\in A$. De plus, puisque $X\in I$, on en déduit que $P\mid X$. Par conséquent, il existe $Q\in A[X]$, de coefficient dominant $q\in A$ tel que $PQ=X$. En passant aux coefficients dominants, on obtient $pq=1$ et $p$ est inversible. Or, de même, il existe $b$ tel que $a=pb$. Or, $P$ est constant et appartient à $I$, donc il existe $c\in A$ tel que $P=p=ac$. On en déduit que $p=pbc$ donc $b$ et $c$ sont inversibles.  Par conséquent, $a=pb$ est inversible comme produit d'inversibles. 
\end{proof}

\begin{ex}[Moulin.AG.61]{Divisibilité et sur-corps}
	Soit $\L$ un sur-corps de $\K$ ainsi que $A$ et $B$ deux polynômes de $\K[X]$. Montrer que $B$ divise $A$ dans $\K[X]$ si, et seulement si, $B$ divise $A$ dans $\L[X]$.
\end{ex}
\begin{proof}
	Écrire les divisions euclidiennes de $A$ par $B$ dans $\L[X]$ et dans $\K[X]$ puis utiliser un argument d'unicité.
\end{proof}

\begin{ex}[Châteaux]{Invariance du PGCD selon le corps de base}
	Soit $\K$ et $\L$ deux corps tels que $\K\subset\L$. Soient $P$ et $Q$ deux polynômes de $\K[X]$.
	\begin{enumerate}
		\item Montrer que le PGCD de $P$ et $Q$ dans $\K[X]$ est le même que dans $\L[X]$.
		\item En déduire que si $P$ et $Q$ sont premiers entre eux dans $\K[X]$ si, et seulement s'ils sont premiers entre eux dans $\L[X]$.
		\item Soit $P$ un polynôme irréductible de $\K[X]$. Montrer que $P$ n'a que des racines simples dans $\L$. Indication : on pourra considérer $P\wedge P'$.
	\end{enumerate}
\end{ex}

\begin{ex}[Moulin.AG.62]{}
	Soit $n\in\N$ et $a\in\R$. Effectuer la division euclidienne dans $\R[X]$ de $$A_n=X^{n+1}\cos((n-1)a)-X^n\cos(na)$$ par $B=X^2-2X\cos(a)+1$.
\end{ex}

\begin{ex}[Moulin.AG.63]{$\star$}
	Pour quelles valeurs de l'entier $n$ le polynôme $(X+1)^n-X^n-1\in\K[X]$ est-il divisible par le polynôme $X^2+X+1$ ? On pourra commencer par $\K=\C$ puis $\K=\R$.
\end{ex}

\begin{ex}[Moulin.AG.64]{Premiers entre eux et racines communes}
	\begin{enumerate}
		\item Montrer que deux polynômes $P$ et $Q$ de $\C[X]$ sont premiers entre eux si, et seulement s'ils n'ont pas de racine complexe commune.
		\item Montrer que deux polynômes $P$ et $Q$ de $\Q[X]$ sont premiers entre eux si, et seulement s'ils n'ont pas de racine complexe commune.
	\end{enumerate}
\end{ex}

\begin{ex}[Moulin.AG.65]{$\star$ Racines de polynômes à coefficients dans $\Q$}
	Soit $P\in\Q[X]$.
	\begin{enumerate}
		\item On suppose $P$ irréductible sur $\Q$. Montrer que les racines complexes de $P$ sont toutes simples. Indication : on pourra utiliser $P'$.
		\item Soit $P\in\Q[X]$ de degré 5. Supposons que $P$ admette une racine double $a$ complexe. Montrer que $P$ possède une racine dans $\Q$.
		\item $\star\star$ On suppose que $P\ne 0$ a une racine $\alpha\in\C$ d'ordre de multiplicité strictement supérieur à $\displaystyle\frac12\deg(P)$. Montrer que $\alpha\in\Q$.
	\end{enumerate}
\end{ex}
\begin{proof}
	\begin{enumerate}
		\item Puisque $P$ est irréductible, $P\wedge P'=1$ puisque cette quantité doit être unitaire, diviser $P$ et $\deg(P')<\deg(P)$. Par conséquent, $X-a$ ne divise pas $P'$ est $a$ est racine simple de $P$ par théorème de cours.
		\item Par contraposée de la question précédente, $P$ n'est pas irréductible et on peut l'écrire sous la forme $P=QR$. SPG, on peut supposer $P$, $Q$ et $R$ unitaires. Si $Q$ ou $R$ est de degré $1$, il n'y a rien à prouver, donc on peut supposer par l'absurde $Q$ et $R$ de degré supérieur à $1$ et irréductibles. Ensuite, puisque $Q\ne R$ comme $\deg(P)=5$, $Q$ et $R$ sont premiers entre eux car leur PGCD vaut $1$, $Q$ ou $R$. Mais si $\deg(Q)=2$, comme $Q$ et $R$ sont irréductibles, d'après l'exercice précédent, on a nécessairement $Q(a)=R(a)=0$ ce qui est absurde. Cela est aussi absurde si $\deg(Q)=3$ car on a alors $\deg(R)=2$.
		\item Raisonner par l'absurde et utiliser la question $1$.
	\end{enumerate}
\end{proof}

\begin{ex}[Moulin.AG.66]{$\Delta$}
	Déterminer les polynômes de $\C[X]$ qui divisent leur polynôme dérivé.
\end{ex}
\begin{proof}
	Utiliser la DES de $P'/P$. Ce sont les $\lambda(X-z)^n$.
\end{proof}

\begin{ex}[Moulin.AG.67]{Composition et divisibilité}
	Soit $(A,B,P)\in\K[X]^3$ avec $\deg(P)\geqslant 1$. Montrer que $B\circ P$ divise $A\circ P$ si, et seulement si, $B$ divise $A$.
\end{ex}

\begin{ex}[Moulin.AG.68]{}
	Soit $P$ et $Q$ deux polynômes de $\R[X]$ premiers entre eux. Montrer que s'il existe $R\in\R[X]$ tel que $P$ divise $R^3-Q$, alors il existe $S\in\R[X]$ tel que $P^2$ divise $S^3-Q$.
\end{ex}

\begin{ex}[Moulin.AG.69]{Résultant}
	\begin{enumerate}
		\item $\Delta$ Soit $P$ et $Q$ dans $\C[X]$ de degrés respectifs $P$ et $Q$. On considère l'application : $$\begin{array}{lrcl}
			\Phi:&\C_{p-1}[X]\times\C_{q-1}[X]&\longrightarrow&\C_{p+q-1}[X]\\
			&(M,N)&\longmapsto&MQ+NP
		\end{array}$$ et l'on pose $\text{Res}(P,Q)$ le déterminant de sa matrice dans les bases canoniques. Montrer que $P$ et $Q$ sont premiers entre eux si, et seulement si, $\text{Res}(P,Q)\ne 0$.
		\item $\star$ Soit $R$ et $S$ dans $\C[X,Y]$ et $\Omega=\left\{(x,y)\in\C^2\ : \ R(x,y)=S(x,y)=0\right\}$. On pose $$\begin{array}{lrcl}
			\pi:&\C^2&\longrightarrow&\C\\
			&(x,y)&\longmapsto&x
		\end{array}$$ Montrer que $\pi(\Omega)$ ou $\C\setminus\pi(\Omega)$ est fini.
	\end{enumerate}
\end{ex}
\begin{proof}
	Pour la première question, utiliser la linéarité de $\Phi$ et déterminer son noyau. On  pourra se référer à la première pâle MP$\st$ 2024-2025.
\end{proof}

\begin{ex}[Moulin.AG.70]{Une amélioration de l'algorithme d'Euclide}
	Montrer que si $A$ et $B$ sont deux polynômes premiers entre eux et non tous deux constants, alors il existe un unique couple $(U,V)\in\K[X]^2$ tel que $AU+BV=1$ avec $\deg(U)<\deg(B)$ et $\deg(V)<\deg(A)$, et que l'algorithme d'Euclide donne ce couple-là (à condition que ...).
\end{ex}
\begin{ex}[Moulin.AG.71]{Presque parties réelles et imaginaires}
	Soit $P$, $Q$ et $R$ trois polynômes de $\C[X]$ tels que $P^2=Q^2+R^2$ et $Q\wedge R=1$. Montrer qu'il existe deux polynômes $P_1$ et $P_2$ premiers entre eux tels que : $$Q=\frac{P_1^2+P_2^2}{2}\ \ \text{et}\ \ R=\frac{P_1^2-P_2^2}{2i}$$
\end{ex}

\begin{ex}[Moulin.AG.72]{$\Delta\star$ Contenu d'un polynôme et applications}
	Soit $P=\displaystyle\sum_{k=0}^{n}a_kX^k\in\Z[X]$. On note $c(P)=\text{PGCD}(a_0,\ldots,a_n)$ le contenu de $P$.
	\begin{enumerate}
		\item Soit $(P,Q)\in\Z[X]^2$. Montrer que $c(PQ)=c(P)c(Q)$. Indication : on pourra se ramener au cas où $c(P)=c(Q)=1$.
		\item Montrer que l'on peut prolonger la fonction $c$ à $\Q[X]$ en conservant la propriété de la première question.
		\item Application 1 : retrouver le fait que les racines rationnelles d'un polynôme unitaire à coefficients entiers sont entières.
		\item Application 2 : soit $P\in\Z[X]$ non constant. Montrer que $P$ est irréductible dans $\Z[X]$ si, et seulement s'il est irréductible dans $\Q[X]$ et $c(P)=1$.
		\item Application 3 : montrer que le polynôme minimal d'une matrice de $\mathcal{M}_n(\Z)$, considérée comme une matrice de $\mathcal{M}_n(\Q)$, est à coefficients entiers.
	\end{enumerate}
\end{ex}
\begin{proof}
	Pour la première question, on procède en deux étapes.
	\begin{itemize}
		\item Supposons $C(A)=C(B)=1$. Par l'absurde, supposons $C(AB)\ne 1$. Prenons alors $p$ un nombre premier qui divise $C(AB)$. Alors, en passant au quotient dans $\Z/p\Z[X]$, on a $\overline{0}=\overline{AB}=\overline{A}\times\overline{B}$. Puisque $\Z/p\Z$ est un corps, $\Z/p\Z[X]$ est intègre donc $\overline{A}=\overline{0}$ ou $\overline{B}=\overline{0}$, et donc $p$ divise $C(A)$ ou $C(B)$, ce qui est absurde. Donc $C(AB)=1$.
		\item Dans le cas général, on écrit $A=C(A)A'$ et $B=C(B)B'$ avec $C(A')=C(B')=1$ (en effet, $C(nP)=\lvert n\rvert C(P)$). Alors, par le cas précédent : $$C(AB)=C\big(C(A)C(B)A'B'\big)=C(A)C(B)C(A'B')=C(A)C(B)$$
	\end{itemize}
\end{proof}

\begin{ex}[Morlot]{$\star$ Polynômes à valeurs dans les nombres premiers}
	On note $\mathcal{P}$ l'ensemble des nombres premiers. Trouver tous les polynômes $A \in \Z[X]$ tels que $\forall n \in \N, \ A(n) \in \mathcal{P}$.
\end{ex}
\begin{proof}
	Utiliser l'astuce classique pour les polynômes $Q$ à coefficients dans $\Z$ : pour tous $a$ et $b$ entiers, $b-a$ divise $Q(b)-Q(a)$. On montre par analyse-synthèse que seuls les polynômes constants égaux à des nombre premiers conviennent.
	\begin{itemize}
		\item \textbf{Analyse} : Soit $A$ qui convienne et on pose $p=A(0)$. Soit $n \in \N\st$. En vertu de la remarque introductive, on sait que $np-0$ divise $A(np)-A(0)=A(np)-p$. Ainsi, on peut fixer $m$ dans $\Z$ tel que $A(np)-p=mnp$. Puis on obtient, $A(np)=(mn+1)p$ donc $p$ divise $A(np)$. Or, $A(np) \in \mathcal{P}$, donc nécessairement, $A(np)=p$. La suite $(np)_{n \in \N}$ étant injective, $A$ et le polynôme constant égal à $p$ coïncident en une infinité de points et sont donc égaux.
		\item \textbf{Synthèse} : Réciproquement, pour tout nombre premier $p$, le polynôme contant égal à $p$ convient.
	\end{itemize}
\end{proof}

\begin{ex}[Morlot]{$\Delta\star$ Critère d'Eisenstein}
	Soit $p$ un nombre premier et $P\in\Z[X]$ \textbf{unitaire} qu'on écrit sous la forme $$P=\sum_{k=0}^{n}a_kX^k$$ Supposons que l'on ait :
	\begin{itemize}
		\item $\forall k\in\llbracket0,n-1\rrbracket,\ p\mid a_k$ ;
		\item $p^2\nmid a_0$.
	\end{itemize}
	Alors $P$ est irréductible dans $\Z[X]$.
\end{ex}
\begin{proof}
	Supposons par l'absurde que $P$ puisse s'écrire sous la forme $P=AB$ avec $A$ et $B$ des polynômes non constants à coefficients entiers. On passe au quotient et on se place dans $\Z/p\Z[X]$. $\Z/p\Z[X]$ est un anneau intègre puisque $\Z/p\Z$ est un corps. Or, on a : $$X^n=\overline{P}=\overline{A}\times\overline{B}$$ puisque $p$ divise les coefficients de $P$ qui ne correspondent pas au terme de degré $N$. Par conséquent, on a $\overline{A}=X^i$ et $\overline{B}=X^j$ avec $i+j=n$ ainsi que $i\geqslant 1$ et $j\geqslant 1$. Ainsi, les coefficients de degré $0$ de $A$ et de $B$ sont divisibles par $p$. Or, le coefficient de degré $0$ de $P$ est égal au produit de celui de $A$ par celui de $B$. Par conséquent, $p^2\mid a_0$, ce qui est absurde.
\end{proof}

\begin{ex}[Morlot]{Une famille de polynômes irréductibles}
	Soit $a_1,\ldots,a_n$ des entiers deux à deux distincts. Montrer que le polynôme $$P=(X-a_1)\times\cdots\times(X-a_n)-1$$ est irréductible sur $\Z[X]$.
\end{ex}
\begin{proof}
	Raisonner par l'absurde en écrivant $P=QR$. On a donc $R(a_i)=\pm 1$ et $Q(a_i)=\mp 1$ (signe opposé). Donc $Q+R$ possède $n$ racines. Discuter ensuite sur les degrés et la limite en $+\infty$ pour conclure à une absurdité.
\end{proof}

\begin{ex}[Révisions]{Polynômes cyclotomiques}
	Soit $n\in\N\st$. On appelle racine primitive $n$-ème de l'unité tout générateur du groupe multiplicatif $\U_n$. On notera $P_n$ l'ensemble des racines primitives $n$-ème de l'unité. Le $n$-ème polynôme cyclotomique est le polynôme : $$\Phi_n=\prod_{x\in P_n}(X-x)$$
	\begin{enumerate}
		\item Quel est le degré de $\Phi_n$ ? Montrer que $$X^n-1=\prod_{d\mid n}\Phi_d$$
		\item Montrer que $\Phi_n$ est à coefficients entiers.
		\item $\star$ Montrer que si $p$ est un nombre premier ne divisant pas $n$, alors $\Phi_n$ est à racines simples sur $\F_p$.
		\item $\star$ Soit $a\in\Z$ et $p$ un nombre premier qui divise $\Phi_n(a)$ mais aucun des $\Phi_d(a)$ pour $d$ diviseur strict de $n$. Montrer que $p\ \equiv\ 1\ [n]$.
		\item Soit $k>n$ et $p$ un diviseur de $\Phi_n(k!)$. Montrer que $p$ est strictement supérieur à $k$ et qu'il est premier avec $n$.
		\item $\star\star$ Montrer que $p\ \equiv\ 1\ [n]$ et en déduire que l'ensemble des nombres premiers congrus à $1$ modulo $n$ est infini.
	\end{enumerate}
\end{ex}
\begin{proof}
	Exercice aberrant mais très joli !
	\begin{enumerate}
		\item $\Phi_n$ est de degré le nombre de générateurs de $\U_n$, soit $\varphi(n)$. Pour la deuxième partie, montrer : $$\left\{\frac{k}{n}, \ k\in\llbracket 1,n\rrbracket\right\}=\bigsqcup_{d\mid n}\left\{\frac{l}{d},\ l\in\llbracket1,d\rrbracket,\ l\wedge d=1\right\}$$ puis scinder $X^n-1$ et regrouper par paquets. \\ Remarque : en passant aux degrés, on retrouve la formule : $$n=\sum_{d\mid n}\varphi(d)$$
		\item Procéder par récurrence forte en utilisant la formule de la première question et en démontrant que le quotient dans la division euclidienne d'un polynôme à coefficients entiers par un polynôme unitaire est un polynôme à coefficients entiers.
		\item On montre que $X^n-1$ est à racines simples dans $\F_p$. En effet, comme $p$ et $n$ sont premiers entre eux, on dispose d'une relation de Bézout : $$ap+bn=1$$ Alors dans $\F_p[X]$: $$bX\times nX^{n-1}-(X^n-1)=bnX^n-(X^n-1)=(1-ap)X^n-(X^n-1)=1$$ si bien que $X^n-1$ et son polynôme dérivé sont premiers entre eux dans $\F_p[X]$ et donc $X^n-1$ est à racines simples dans $\F_p[X]$. Comme $\phi_n$ divise $X^n-1$, $\Phi_n$ est aussi à racines simples dans $\F_p$.
		\item $a^n=1$ dans $\F_p$ si bien que l'ordre de $a$ dans $\F_p\st$ divise $n$. Notons $o$ l'ordre de $a$ dans $\F_p\st$. Alors d'après la question $1$, on a dans $\F_p$ : $$0=a^o-1=\prod_{d\mid o}\Phi_d(a)$$ Or, le seul facteur potentiellement nul dans le membre de droite correspond à $d=n$, donc $n\mid o$ et $o=n$. D'après le théorème de Lagrange, on en déduit que $n$ divise $p-1$ et donc le résultat.
		\item D'après la question $1$, $(k!)^n-1$ est divisible par $p$ donc $p$ est premier avec $(k!)^n$, puis $p$ est premier avec $k!$ car $p$ est premier, et ainsi $p>k$. Donc $p>k>n$ et comme $p$ est premier, $p\wedge n=1$.
		\item De $5$ et du lemme démontré en $3$, on en déduit que les $\Phi_d(a)$ pour $d$ un diviseur de $n$ différent de $n$ ne sont pas divisibles par $p$. Par $4$, $p$ est congru à $1$ modulo $n$. \\ Pour conclure, on prend $p_1$ le nombre premier obtenu en prenant $k=n+1$, puis $p_2$ celui obtenu en prenant $k=p_1+1$, etc. On obtient ainsi une suite strictement croissante de nombres premiers congrus à $1$ modulo $n$, donc une infinité !
	\end{enumerate}
\end{proof}

\end{document}